\begin{frame}<1-5>[t]{新知探索}
\begin{campuscanvas}
% Source slide 8, objects 5; visible 1-5
\only<1-5|handout:1>{\node[anchor=north west,inner sep=0,text=black] at (70.000,150.000) {\begin{minipage}[t]{142.500mm}\raggedright\fontsize{14}{21.00}\selectfont \textbf{思考}\quad 军训时教官喊 13 班集合：\end{minipage}};}
% Source slide 8, objects 8; visible 2-5
\only<2-5|handout:1>{\node[anchor=north west,inner sep=0,text=black] at (70.000,255.000) {\begin{minipage}[t]{100.000mm}\raggedright\fontsize{12}{18.00}\selectfont 14 班学生会不会跑到 13 班来？\end{minipage}};}
% Source slide 8, objects 8; visible 2-5
\only<2-5|handout:1>{\node[anchor=north west,inner sep=0,text=black] at (70.000,360.000) {\begin{minipage}[t]{100.000mm}\raggedright\fontsize{12}{18.00}\selectfont 教官调整了站位后班级里的人有没有发生变化？\par 班级会不会发生改变？\end{minipage}};}
% Source slide 8, objects 8; visible 2-5
\only<2-5|handout:1>{\node[anchor=north west,inner sep=0,text=black] at (70.000,535.000) {\begin{minipage}[t]{100.000mm}\raggedright\fontsize{12}{18.00}\selectfont 教官要求报数的目的是什么？一个人是否会报两次？\end{minipage}};}
% Source slide 8, objects 9; visible 3-5
\only<3-5|handout:1>{\node[anchor=north west,inner sep=0,text=black] at (970.000,255.000) {\begin{minipage}[t]{28.750mm}\raggedright\fontsize{14}{21.00}\selectfont \red{确定性}\end{minipage}};}
% Source slide 8, objects 11; visible 4-5
\only<4-5|handout:1>{\node[anchor=north west,inner sep=0,text=black] at (970.000,380.000) {\begin{minipage}[t]{28.750mm}\raggedright\fontsize{14}{21.00}\selectfont \red{无序性}\end{minipage}};}
% Source slide 8, objects 12; visible 5
\only<5|handout:1>{\node[anchor=north west,inner sep=0,text=black] at (970.000,545.000) {\begin{minipage}[t]{28.750mm}\raggedright\fontsize{14}{21.00}\selectfont \red{互异性}\end{minipage}};}
\end{campuscanvas}
\end{frame}
